\section{A single Pipeline}
\label{sec:pipeline}

All the preprocessing and the model are wrapped in one scikit-learn
\code{Pipeline}, which takes a raw census row and returns $P(\pos)$:
(1) a stateless \code{FunctionTransformer} adds the three engineered features;
(2) a \code{ColumnTransformer} applies median imputation and
\code{StandardScaler} to the numerical columns, most-frequent imputation and
\code{OneHotEncoder(handle\_unknown="ignore")} to the categorical ones, and
drops \code{education} and \code{fnlwgt} (\code{remainder="drop"}); (3) the
classifier. One Pipeline brings four guarantees.
\begin{itemize}
  \item \textbf{No leakage during cross-validation.} A fresh, unfitted copy
  (\code{clone}) is fitted in each fold, so the medians, means, standard
  deviations and categories are learned on the four training folds only.
  Imputing or scaling the whole development set first would have let each
  validation fold influence its own preprocessing.
  \item \textbf{The same transformations everywhere.} The same object is
  cross-validated, refitted on the whole development set, evaluated on the
  test set and saved: no second implementation of the preprocessing can drift
  from the first one.
  \item \textbf{Tuning covers the whole chain.} Each Optuna trial builds and
  cross-validates a complete Pipeline (Section~\ref{sec:tuning}).
  \item \textbf{One file to deploy.} The fitted Pipeline is saved with
  \code{joblib}, together with its decision threshold, since both are needed
  to turn a raw row into a decision. It also returns the name of each encoded
  column (\code{get\_feature\_names\_out}), which the SHAP analysis uses.
\end{itemize}
